\section{O que é o Docker?}

\begin{frame}
  \frametitle{O Problema}
  Uma aplicação com três componentes:

  \vspace{0.5em}
  \begin{itemize}
    \item \textbf{React frontend} — precisa de Node
    \item \textbf{Python API} — precisa de Python
    \item \textbf{PostgreSQL} — precisa do banco instalado
  \end{itemize}

  \vspace{0.5em}
  Como garantir que as versões sejam as mesmas em todo o time, no CI/CD
  e sem interferência de outras ferramentas locais?
\end{frame}

\begin{frame}
  \frametitle{A Solução: Containers}
  Um \textbf{container} é um processo isolado para cada componente da aplicação.

  \vspace{0.5em}
  Cada parte roda em seu próprio ambiente — completamente separado da máquina local.

  \vspace{0.5em}
  \begin{itemize}
    \item \textbf{Self-contained} — tudo que precisa, sem depender de dependências locais
    \item \textbf{Isolado} — influência mínima da máquina e de outros containers
    \item \textbf{Independente} — deletar um container não afeta os outros
    \item \textbf{Portável} — roda igual em qualquer máquina, servidor ou nuvem
  \end{itemize}
\end{frame}

\begin{frame}
  \frametitle{Containers vs Máquinas Virtuais}
  \begin{columns}
    \column{0.5\textwidth}
      \textbf{Máquina Virtual}
      \begin{itemize}
        \item Sistema operacional completo
        \item Kernel, drivers e programas próprios
        \item Pesada e lenta para iniciar
        \item Pouco eficiente para isolar uma aplicação
      \end{itemize}

    \column{0.5\textwidth}
      \textbf{Container}
      \begin{itemize}
        \item Processo isolado e leve
        \item Compartilha o kernel do \textit{host}
        \item Rápido para iniciar
        \item Mais aplicações com menos infraestrutura
      \end{itemize}
  \end{columns}
\end{frame}